\small
\setlength{\parskip}{0pt}

\vkrformministry

\vspace{0.4em}
\vkrformapprove

\vspace{0.5em}
\begin{center}
{\bfseries\Large З~А~Д~А~Н~И~Е}\\[0.25em]
{\bfseries\large на выполнение выпускной квалификационной работы}
\end{center}

\vspace{0.4em}
\noindent Студент группы \underline{\hspace{0.3em}ИУ9-82Б\hspace{0.3em}}

\vspace{0.2em}
\noindent\underline{\makebox[0.92\textwidth]{Аринин Матвей Алексеевич}}\\
\hbox to \textwidth{\hfil\scriptsize{(фамилия, имя, отчество)}\hfil}

\vspace{0.4em}
\noindent Тема квалификационной работы: \uline{\hspace{0.35em}Компилятор базисного подмножества Рефала-5 с представлением данных, использующим раскрученные очереди}
\vspace{0.1em}

\noindent\underline{\hfill}

\vspace{0.4em}
\noindent При выполнении ВКР:

\vspace{0.15em}
\noindent
\begin{tabular}{|p{\dimexpr\textwidth-2.05cm\relax}|c|}
\hline
\centering Используются / Не используются & Да/Нет \\
\hline
1) Литературные источники и документы, имеющие гриф секретности & Нет \\
\hline
\begin{tabular}[t]{@{}p{\linewidth}@{}}
2) Литературные источники и документы, имеющие пометку <<Для служебного пользования>>, иных пометок, запрещающих открытое опубликование
\end{tabular} & Нет \\
\hline
3) Служебные материалы других организаций & Нет \\
\hline
4) Результаты НИР (ОКР), выполняемой в МГТУ им. Н.Э. Баумана & Нет \\
\hline
\begin{tabular}[t]{@{}p{\linewidth}@{}}
5) Материалы по незавершённым исследованиям или материалы по завершённым исследованиям, но ещё не опубликованные в открытой печати
\end{tabular} & Нет \\
\hline
\end{tabular}

\vspace{0.6em}
\noindent\makebox[\textwidth][s]{Тема квалификационной работы утверждена распоряжением по факультету}

\vspace{0.3em}
\noindent\underline{\hspace{4.0cm}} № \underline{\hspace{2.8cm}} от <<\underline{\hspace{0.8cm}}>> \underline{\hspace{2.8cm}} 20\underline{\hspace{0.8cm}} г.

\vspace{0.5em}
\begingroup
\setlength{\parindent}{1.25cm}
\setlength{\parskip}{0.7em}

\textbf{Часть 1.} Провести обзор и сравнительный анализ распространённых представлений объектных выражений в реализациях Рефала-5 и обосновать выбор раскрученной двусторонней очереди.

\newpage
\textbf{Часть 2.} Разработать архитектуру компилятора базисного подмножества Рефала-5, включая состав модулей трансляции и исполнения, интерфейсы и форматы внутренних данных.

\textbf{Часть 3.} Выполнить программную реализацию компилятора, провести автоматическое тестирование и сравнительный анализ производительности с альтернативными реализациями Рефала-5.

\noindent{\bfseries\itshape Оформление квалификационной работы:}

\noindent Расчётно-пояснительная записка на \underline{\hspace{1.8cm}} листах формата А4.

\noindent Перечень графического (иллюстративного) материала (чертежи, плакаты, слайды и т.п.)\\
\underline{Презентация\hfill}

\noindent Дата выдачи задания <<\underline{\hspace{0.7cm}}>> \underline{\hspace{2.4cm}} 20\underline{\hspace{0.7cm}} г.

В соответствии с учебным планом выпускную квалификационную работу выполнить в полном объёме в срок до <<\underline{\hspace{0.7cm}}>> \underline{\hspace{2.4cm}} 20\underline{\hspace{0.7cm}} г.
\endgroup

\vspace{1.0em}
\newcommand{\vkrzadsig}[2]{%
\noindent\makebox[\dimexpr\textwidth-7.6cm\relax][l]{\bfseries #1}\underline{\makebox[3.6cm]{\strut}}\hspace{0.4cm}\underline{\makebox[3.6cm]{\strut #2}}\\[-0.4ex]
\noindent\makebox[\dimexpr\textwidth-7.6cm\relax][l]{}\makebox[3.6cm][c]{\scriptsize(Подпись, дата)}\hspace{0.4cm}\makebox[3.6cm][c]{\scriptsize(И.О. Фамилия)}\\[0.8em]}
\vkrzadsig{Руководитель квалификационной работы}{}
\vkrzadsig{Студент}{М.А. Аринин}

\vspace{0.8em}
\noindent\uline{Примечание}:\\
1. Задание оформляется в двух экземплярах: один выдаётся студенту, второй хранится на кафедре.
